\scanpage{59}
$E$ \uline{catégorie}. Conditions équivalentes

a) \deleted{Il existe une} $E$ est stable par petites $\varinjlim$ et admet une petite \inserted{sous-}catégorie pleine str génératrice

a$'$) $E$ est stable par petites $\varinjlim$, et il existe une petite catégorie $C$ et un foncteur $\varphi:C\mathbin{-}E$ str. générateur

a$''$) Il existe une petite catégorie $C$, et un foncteur \uline{pleinement fidèle} $E\to\widehat C$ qui admet un adjoint à gauche $\widehat C\mathbin{-}E$.

[Dans une telle catégorie, tt \inserted{contra}foncteur \deleted{\unsure{en}} dans (Ens) qui commute aux petites $\varprojlim$ est représentable, donc il existe ds $E$ des petites $\varprojlim$, et dans a$''$), le foncteur $E\mathbin{-}\widehat C$ y commute, i.e. $E$ s'identifie à une sous-catégorie stable par $\varprojlim$ de $\widehat C$. Si $E$ est accessible [ce qui est plus fort que les conditions équivalentes a)\ldots], alors : \deleted{\unsure{Il y a un}} un foncteur $E\to\widehat C$ a un adjoint à g ssi il est \uline{accessible} et commute aux \inserted{petites} $\varprojlim$]

\deleted{\unsure{Th}} \uline{Question} Une catégorie satisfaisant les conditions a) est-elle \uline{accessible}, i.e. ses éléments sont ils accessibles ? Comme, sous les conditions de a$''$), $\widehat C$ est accessible, on voit qu'il \inserted{faut et} suffit que le foncteur $E\to\widehat C$ soit accessible, \deleted{ou} ce qui revient au \unsure{m\textsuperscript{e}}, qu'il existe un cardinal $\pi$ tel que $E$ soit stable dans $\widehat C$ par $\varinjlim$ filtrantes grandes devant $\pi$.

\margin{Contre-exemple : $(\mathrm{Ens})^\circ$, elle \deleted{\unsure{admet}} a des \inserted{petites} limites inductives et un cogénérateur strict (savoir tt ens de cardinal $\geq 2$) mais \deleted{\unsure{en fait}} les seuls objets accessibles sont les ensembles ponctuels !}% slanted

\uline{Exemples} a) Topos \deleted{\unsure{a) opposée d'une catégorie}} \inserted{a') \unsure{opposée|opposés|opposé} de topos.}

b) $T(\mathrm{Ens})$, $T$ une structure algébrique définissable par petites $\varprojlim$

c) $T(C)$, où $T$ est comme dans b), et $C\deleted{\unsure{\}}}$ une catégorie \uline{accessible}

\deleted{\unsure{ensembliste}}

\margin{Conséquence : 1) Si $T\mathbin{-}T'$ est un hom de structures algébriques, d'où $T(C)\to T'(C)$ foncteur \uline{accessible} entre catégories accessibles, commutant aux $\varprojlim$ petites, donc $\exists$ adjoint à gauche (structures $T'$-$T$-libres) 2) Si $C\to C'$ foncteur accessible commutant aux pet. $\varprojlim$ entre cat. acc. $T(C)\to T(C')$ est \uline{accessible} commutant aux petites $\varprojlim$, donc $\exists$ adjoint à g\ldots}

À prouver : Si $C,C'$ accessibles, alors la ss-catégorie de $\underline{\operatorname{Hom}}(C^\circ\times C'^\circ,\mathrm{Ens})$ \inserted{$\underline{\operatorname{Hom}}^*(C^\circ\times C'^\circ,\mathrm{Ens})$} formée des bifoncteurs qui commutent aux $\varprojlim$ \unsure{quelc.|quelq.|petites} en chaque argument, $(\simeq\underline{\operatorname{Hom}}'(C^\circ,C')\simeq\underline{\operatorname{Hom}}'(C'^\circ,C)$, où les $\underline{\operatorname{Hom}}'$ dénotent les foncteurs qui commutent aux petites $\varprojlim$) est accessible, et des foncteurs commutant aux $\varprojlim$ $C_1\mathbin{-}C,C'_1\mathbin{-}C'$ définissent un foncteur accessible entre les $\underline{\operatorname{Hom}}^*$. Or il est assez clair que plus généralement, si $D$ \inserted{variable} stable par petites $\varprojlim$ (et pour les foncteurs qui y commutent\ldots cf \inserted{note} $\Delta$)
\[
 \underline{\operatorname{Hom}}^*(C^\circ\times C'^\circ,\unsure{D|\mathcal D})\simeq T_{C^\circ}T_{C'^\circ}(D)\simeq T_{C'^\circ}T_{C^\circ}(D)
\]
est une ``théorie algébrique'', définie par $(C*C')^\circ$ [$C*C'$ universel pour ``accouplements'' $C\times C'\mathbin{-}C''$ commutant aux $\varinjlim$ en chaque argument\ldots]
% archive sheet number: 58
